\section{从打开文章到内容加载}
\subsection{一个页面其实有几种不同的加载}
在文章页任务中，我开始区分代码、数据与媒体资源。代码决定页面怎么运行，接口数据提供文章标题和正文，%
图片文件提供实际画面。项目的路由动态 import 解决页面代码什么时候下载的问题，%
接口请求解决内容什么时候到达的问题，图片懒加载解决哪些图片暂时不需要读取的问题。%
它们都叫“按需加载”，却作用在不同对象上。

Vue Router 支持以动态导入定义路由组件，浏览器访问对应路由时再获取相关代码\cite{router-lazy}。%
这适合把后台编辑器等非首屏功能从普通访客的初始路径中分离出来。%
代价是第一次切换到该页面时可能增加一次等待，所以仍需要清楚的加载反馈，%
并检查共享依赖是否被其他入口提前引入。写了动态 import 并不能保证所有资源都已实现最优拆分。

数据加载则需要考虑请求之间的关系。当前文章详情先读取正文，关联商品读取失败时仍保留文章；%
列表使用请求序号，避免较早搜索的慢响应覆盖新的筛选结果\cite{local-code}。%
我把它理解为“只接纳最新一次问题的答案”，而不是“谁最后送到就显示谁”。取消请求可以减少无效工作，%
但提交结果前仍要检查它是否属于当前页面意图。

\subsection{正文在浏览器生成，还是提前生成}
客户端渲染（CSR）主要在浏览器执行 JavaScript 后形成业务正文；服务端渲染（SSR）在处理请求时生成
HTML；静态生成（SSG）在构建时提前生成 HTML\cite{vue-ssr}。水合则是客户端把组件状态和事件连接到已存在
HTML 的过程。对初学者而言，最容易抓住的区别是：用户拿到初始 HTML 时，%
里面是否已经有可阅读的文章内容。

\begin{figure}[H]
\centering
\begin{tikzpicture}[
  font=\fontsize{10.5}{14}\selectfont,
  stage/.style={draw=wmnavy!65,fill=white,rounded corners=2pt,
    text width=3.0cm,minimum height=1.35cm,align=center,
    inner sep=4pt,line width=0.6pt},
  link/.style={-{Stealth[length=2mm]},draw=wmaccent,line width=0.7pt}
]
\node[anchor=west,font=\fontsize{11}{14}\selectfont\bfseries,text=wmnavy]
  at (0,1.05) {CSR：客户端渲染};
\node[stage] (c1) at (1.65,0) {浏览器\\请求页面};
\node[stage,fill=wmnavy!5] (c2) at (5.4,0) {返回 HTML 壳\\通常缺少主要正文};
\node[stage] (c3) at (9.15,0) {浏览器执行 JS\\获取数据并渲染};
\node[stage,fill=wmnavy!9] (c4) at (12.9,0) {主要正文出现\\建立页面交互};
\draw[link] (c1) -- (c2);
\draw[link] (c2) -- (c3);
\draw[link] (c3) -- (c4);
\node[anchor=west,font=\fontsize{11}{14}\selectfont\bfseries,text=wmnavy]
  at (0,-1.35) {SSR：服务端渲染};
\node[stage] (r1) at (1.65,-2.4) {浏览器\\请求页面};
\node[stage,fill=wmnavy!9] (r2) at (5.4,-2.4) {服务端生成页面\\HTML 已含正文};
\node[stage,fill=wmnavy!9] (r3) at (9.15,-2.4) {浏览器显示正文\\此时可未执行 JS};
\node[stage] (r4) at (12.9,-2.4) {执行 JS 并水合\\接管已有界面交互};
\draw[link] (r1) -- (r2);
\draw[link] (r2) -- (r3);
\draw[link] (r3) -- (r4);
\node[anchor=west,font=\fontsize{11}{14}\selectfont\bfseries,text=wmnavy]
  at (0,-3.75) {SSG：静态生成};
\node[stage] (s1) at (1.65,-4.8) {构建时生成\\含正文的 HTML};
\node[stage,fill=wmnavy!9] (s2) at (5.4,-4.8) {页面请求到达\\返回静态 HTML};
\node[stage,fill=wmnavy!9] (s3) at (9.15,-4.8) {浏览器显示正文\\此时可未执行 JS};
\node[stage] (s4) at (12.9,-4.8) {按需执行 JS\\增强交互（可选）};
\draw[link,dashed] (s1) -- (s2);
\draw[link] (s2) -- (s3);
\draw[link] (s3) -- (s4);
\node[anchor=north west,text width=14.45cm,align=left,text=wmnavy]
  at (0,-5.85)
  {典型流程示意：底色强调正文何时可用；箭头不表示耗时比例。SSG 的构建先于页面请求；%
    混合渲染可按路由选择策略。};
\end{tikzpicture}

\caption{理解三种渲染方式：正文在哪一步出现}%
\label{fig:render-path}
\end{figure}

当前项目采用 CSR，文章数据和页面标题由客户端更新。它便于与编辑、筛选等交互结合，%
也符合现有共享应用的组织方式；缺点是正文展示依赖代码执行与数据请求，初始 HTML 的文章信息有限。%
SSR 可以更早交付正文，但引入服务器运行、状态隔离和水合一致性问题；SSG 的静态交付简单，%
却需要在内容发布或下线时及时更新产物。不存在一种方案在交互、维护成本和更新时效上同时没有代价。

现阶段的发展并非只有“全部 CSR”或“全部 SSR”两种选择。Nuxt 提供按路由区分渲染与缓存策略的能力；%
Astro 的岛屿架构则把少量交互组件放在以静态内容为主的页面中\cite{nuxt-rendering,astro-islands}。%
其思想类似一篇文章只有筛选器需要复杂交互，就不必让整篇正文都承担同样的客户端工作。%
这类方案值得关注，但并不要求本项目立即整体迁移。


对 WEMOVE，我认为后台编辑继续使用客户端交互是合理的；如果公开文章以后更重视搜索和分享预览，%
可以选择一条文章路由试点预渲染或 SSR。迁移时必须考虑正文目前使用 DOMParser 等浏览器 API，%
以及文章下线后旧 HTML 如何立即失效。公开文章也不意味着整页都可以公共缓存，%
导航中的昵称和管理员入口仍要独立处理。

\subsection{图片、分页和虚拟列表解决不同问题}
图片常常比文字和页面代码占用更多字节。项目既有静态统计显示，六张 $1254\times1254$ 的玩法 PNG 合计约
9.13 MiB；这不代表一个页面会同时下载六张，却足以提示我关注图片规格。把大图放进较小容器，CSS
只改变显示尺寸，并不会自动把下载文件变小。

\begin{figure}[H]
\centering
\includegraphics[width=0.68\linewidth]{online/webdev-responsive-image-layout.png}
\caption{图片超出视口与适应视口的直观对比}%
\label{fig:loading-online}
\end{figure}

图\ref{fig:loading-online}说明了图片首先应在页面中完整、合理地显示。%
但显示合适与下载合适是两个问题，后者还需要资源规格配合。

响应式图片通过 srcset 提供多个候选资源，通过 sizes 描述显示条件，%
让浏览器结合布局与像素密度选择图片\cite{responsive-images}。优点是减少不必要的大图传输，%
代价是需要生成和维护多尺寸文件。当前正文图片有懒加载和异步解码，但安全规则只允许受控地址，%
不能直接让文章作者写入任意 srcset。合理的扩展应由媒体服务提供可信候选，而不是绕过白名单。

\begin{table}[!htb]
\centering
\smalltable
\caption{容易被混为一谈的加载优化}
\begin{tabularx}{\linewidth}{>{\raggedright\arraybackslash}p{2.45cm}YY}
\toprule
技术 & 主要减少什么 & 使用时的注意点\\
\midrule
路由懒加载 & 尚未访问页面的初始代码 & 首次导航存在等待，应提供状态反馈\\
\tabledivider
图片懒加载 & 暂时不需要显示的图片请求 & 真正的首屏关键图片不宜延迟发现\\
\tabledivider
分页 & 单次读取和显示的数据数量 & 需要维护页码、筛选条件与返回行为\\
\tabledivider
虚拟列表 & 长列表同时存在的 DOM 节点 & 高度、焦点与滚动恢复更复杂，不会自动减少全部数据请求\\
\tabledivider
流式内容交付 & 等待完整结果后才开始显示的时间 & 需处理分段错误、加载边界和渐进呈现\\
\bottomrule
\end{tabularx}
\end{table}

由此我认识到，“越晚加载越快”也不准确。首屏最重要的图片应尽早被发现，非首屏图片再按需获取；%
用户主动打开编辑页后，必需的编辑能力也不应被无意义地拖延。性能优化更像安排工作顺序：%
先完成影响当前任务的部分，再处理其余内容，而不是对所有资源使用同一个开关。
